\subsection{\texorpdfstring{SYSTEMS OF TYPE \(C_{l}\) (\(l \geqslant 2\))}{SYSTEMS OF TYPE C\_l (l >= 2)}}

\begin{enumerate}
\def\labelenumi{(\Roman{enumi})}
\item
  The existence of root systems of type \(C_{l}\) has been proved in no. 5, since we have seen that the inverse system of a system of type \(B_{l}\) is of type \(C_{l}\) . A root system of type \(C_{l}\) is thus obtained by taking in \(R^{l}\) the vectors \(\pm2\varepsilon_{i}\ (1\leqslant i\leqslant l)\) , and \(\pm\varepsilon_{i}\pm\varepsilon_{j}\ (1\leqslant i<j\leqslant l)\) . The number of roots is \(2l^{2}\) .
\item
  A basis of R can be obtained by taking the image under the map \(\alpha\mapsto\frac{2\alpha}{(\alpha|\alpha)}\) of the basis of the system considered in no. 5. We obtain:
\end{enumerate}

\[
\alpha_ {1} = \varepsilon_ {1} - \varepsilon_ {2}, \alpha_ {2} = \varepsilon_ {2} - \varepsilon_ {3}, \dots , \alpha_ {l - 1} = \varepsilon_ {l - 1} - \varepsilon_ {l}, \alpha_ {l} = 2 \varepsilon_ {l}.
\]

The positive roots are the \(2\varepsilon_{i}\) and the \(\varepsilon_{i} \pm \varepsilon_{j} (i < j)\) .

\begin{enumerate}
\def\labelenumi{(\Roman{enumi})}
\setcounter{enumi}{2}
\item
  The Coxeter number is the same as for the inverse system: \(h = 2l\) .
\item
  Let \(\tilde{\alpha}=2\varepsilon_{1}=2\alpha_{1}+2\alpha_{2}+\cdots+2\alpha_{l-1}+\alpha_{l}\) , which is a root. The sum of its coordinates relative to \((\alpha_{i})\) is 2l-1=h-1. Hence, \(\tilde{\alpha}\) is the highest root. We have \((\tilde{\alpha}|\alpha_{i})=0\) for \(i\neq l\) , \((\tilde{\alpha}|\alpha_{l})=2\) . Thus, the completed Dynkin graph is
\end{enumerate}

\[
\alpha_ {1} \quad \alpha_ {2} \quad \dots \quad \alpha_ {l - 2} \quad \alpha_ {l - 1} \quad \alpha_ {l}
\]

\begin{enumerate}
\def\labelenumi{(\Alph{enumi})}
\setcounter{enumi}{21}
\tightlist
\item
  We have already determined \(R^{-}\) , which is of type \(B_{l}\) . By formula (19) of § 1, no. 12 and by no. 5 (V), the square of the length of \(2\varepsilon_{i}\) for \(\Phi_{R}\) is
\end{enumerate}

\[
((l + 1) (4 l - 2)) ^ {+ 1} ((4 l - 2) ^ {- 1}) ^ {- 1} = (l + 1) ^ {- 1};
\]

thus, \(\Phi_{\mathrm{R}}(x,y) = (x|y) / 4(l + 1)\) .

We have \(\gamma(R) = \gamma(R') = (l + 1)(4l - 2)\) .

\begin{enumerate}
\def\labelenumi{(\Roman{enumi})}
\setcounter{enumi}{5}
\tightlist
\item
  The fundamental weights are easily found:
\end{enumerate}

\[
\begin{array}{l} \omega_ {i} = \varepsilon_ {1} + \varepsilon_ {2} + \dots + \varepsilon_ {i} \\ = \alpha_ {1} + 2 \alpha_ {2} + \dots + (i - 1) \alpha_ {i - 1} + i (\alpha_ {i} + \alpha_ {i + 1} + \dots + \frac {1}{2} \alpha_ {l}) (i \leqslant l). \end{array}
\]

\begin{enumerate}
\def\labelenumi{(\Roman{enumi})}
\setcounter{enumi}{6}
\tightlist
\item
  The sum of the positive roots is
\end{enumerate}

\[
\begin{array}{r l} 2 \rho & = 2 l \varepsilon_ {1} + (2 l - 2) \varepsilon_ {2} + \dots + 4 \varepsilon_ {l - 1} + 2 \varepsilon_ {l} \\ & = 2 l \alpha_ {1} + 2 (2 l - 1) \alpha_ {3} + \dots + i (2 l - i + 1) \alpha_ {i} + \dots \\ & \qquad \dots + (l - 1) (l + 2) \alpha_ {l - 1} + \frac {1}{2} l (l + 1) \alpha_ {l}. \end{array}
\]

\begin{enumerate}
\def\labelenumi{(\Roman{enumi})}
\setcounter{enumi}{7}
\item
  By no. 4 and no. 5 (VIII), \(Q(R) = L_{1}\) , \(P(R) = L_{0}\) ; \(P(R)/Q(R)\) is isomorphic to Z/2Z, and the connection index is 2.
\item
  and (X) These data depend only on W(R), and so are the same as for type \(B_{l}\) .
\item
  The same argument as in no. 5 shows that \(A(R) = W(R)\) and \(w_0 = -1\) .
\item
  The single non-identity element of \(\mathrm{P}(\mathrm{R}^{-})/\mathrm{Q}(\mathrm{R}^{-})\) defines the unique non-trivial automorphism of the completed Dynkin graph: it interchanges the vertices corresponding to \(\alpha_{j}\) and \(\alpha_{l-j}\) for \(0 \leqslant j \leqslant l\) .
\end{enumerate}

